\documentclass[conference]{IEEEtran}
\IEEEoverridecommandlockouts
\usepackage{cite}
\usepackage{amsmath,amssymb,amsfonts,amsthm}
\usepackage{algorithmic}
\usepackage{graphicx}
\usepackage{textcomp}
\usepackage{xcolor}
\usepackage{booktabs}
\usepackage{url}

\newtheorem{theorem}{Theorem}
\newtheorem{lemma}{Lemma}

\begin{document}

\title{Privacy-Preserving Biometric Template Protection via Five-Impression Consensus Fuzzy Commitment and Ristretto255 Schnorr Zero-Knowledge Proofs}

\author{\IEEEauthorblockN{Himakshi Sharma}
\IEEEauthorblockA{\textit{Department of Computer Science and Engineering} \\
\textit{SRM University-AP}\\
Andhra Pradesh, India \\
himakshi\_sharma@srmap.edu.in}
}

\maketitle

\begin{abstract}
Traditional biometric authentication systems store raw ridge images or unencrypted minutiae vectors in centralized databases, leaving user biological identities irrevocably compromised upon database exfiltration. While cryptographic Fuzzy Commitment Schemes (FCS) protect templates by binding them to random keys, they typically exhibit two critical flaws: (1) the recovered secret or its hash is transmitted directly to the server during authentication, defeating end-to-end privacy, and (2) intra-class biological variations cause high False Rejection Rates (FRR), inflating the Equal Error Rate (EER) well above operational thresholds. Conversely, recent general-purpose zero-knowledge proof (zk-SNARK) biometric schemes suffer from prohibitive prover times (3,000--25,000 ms) and excessive memory footprints. In this paper, we propose an ultra-lightweight, end-to-end privacy-preserving biometric authentication architecture that resolves this trilemma. The framework integrates an audited 3-channel 384-bit annular orientation encoder, an algebraic Fuzzy Commitment Scheme with interleaved majority-voting Block Error-Correcting Codes (Cfg-32: $N=384, K=32, R=12$), a five-impression consensus enrollment binding protocol with rotational alignment and reliability-weighted template fusion, a two-factor HKDF witness derivation, and an ultra-fast 3-pass Schnorr Sigma Zero-Knowledge Proof protocol over the prime-order Ristretto255 group. Evaluated across benchmark multi-sample fingerprint datasets, five-impression consensus fusion expands active annular template coverage from 91.4\% to 96.1\%, increases cross-impression template overlap to 88.4\%, and boosts exact biometric secret recovery by +18.0 percentage points compared to single-image enrollment, achieving an empirical 0.00\% cryptographic False Acceptance Rate (FAR) under random-token attacks. Furthermore, zero-knowledge proof generation and verification execute in just 5.47 ms with a communication payload of 160 bytes, achieving over a $500\times$ speedup compared to circuit-based zk-SNARKs with rigorous ISO/IEC 30136 unlinkability ($D_{sys} \to 0$).
\end{abstract}

\begin{IEEEkeywords}
Biometric Template Protection, Fuzzy Commitment, Zero-Knowledge Proofs, Ristretto255, Schnorr Sigma Protocol, Multi-Impression Fusion, ISO/IEC 30136 Unlinkability.
\end{IEEEkeywords}

\section{Introduction}
Biometric authentication has largely replaced alphanumeric credentials in consumer mobile devices, identity verification, and access control. However, biometrics possess an immutable property: unlike passwords or cryptographic keys, human physical features cannot be revoked or re-issued once compromised. Centralized storage of unencrypted biometric templates presents an existential security risk.

Existing attempts to secure biometric templates face an architectural trilemma:
\begin{enumerate}
    \item \textbf{The Server-Leakage Flaw:} Classic Fuzzy Commitment schemes~\cite{juels1999fuzzy} send the candidate secret $S'$ or its hash to the verifier upon login. An untrusted server holding both helper data $W$ and $S'$ can compute $B = W \oplus \text{Codeword}(S')$, instantly reconstructing the user's unencrypted biological template.
    \item \textbf{The Computational Bottleneck:} Recent zk-SNARK biometric schemes~\cite{zkbiometrics2024} convert distance calculations into arithmetic circuits. Generating proofs for million-gate circuits takes 3.5 to 25 seconds on mobile CPUs, rendering them unusable in interactive environments.
    \item \textbf{Intra-Class Error Degradation:} Biological fingerprints suffer from non-linear skin elasticity, rotational misalignment, and sensor moisture. Cryptographic codes that cannot absorb these spatial variations suffer high False Rejection Rates (FRR), inflating the Equal Error Rate (EER).
\end{enumerate}

To address these challenges, we propose an end-to-end architecture that achieves \textbf{robust error tolerance} and \textbf{true zero server leakage} with \textbf{millisecond-level execution}.

\subsection{Novel Contributions}
\begin{itemize}
    \item \textbf{Five-Impression Consensus Template Fusion:} We design a mask-aware, reliability-weighted fusion algorithm that aligns five enrollment impressions across 16 cyclic angular shifts, selects a high-coverage reference impression, and synthesizes a single stable 384-bit orientation template with 96.1\% active annular coverage.
    \item \textbf{Zero Server-Side Leakage via Ristretto255 Schnorr ZKP:} The server database stores solely cryptographic noise: Helper Data $W$, Mask $M$, and Ristretto255 Public Key $Y$. The user proves knowledge of an HKDF-derived secret witness $x = \text{HKDF}(S, \text{salt}=\text{uid}, \text{info}=K)$ without revealing $S$, $K$, minutiae, or ridge images.
    \item \textbf{Millisecond Prover Latency ($<$6 ms):} Operating over the prime-order Ristretto255 group (RFC 9496), proof generation executes in 3.05 ms and verification in 2.42 ms (5.47 ms total), with a transmission footprint of just 160 bytes.
    \item \textbf{Formal Security and Privacy Proofs:} We mathematically prove information-theoretic helper secrecy, special soundness, honest-verifier zero-knowledge (HVZK), and ISO/IEC 30136 unlinkability ($D_{sys} \to 0$).
\end{itemize}

\section{Proposed Architecture}

\subsection{3-Channel Sector-Quantized Annular Feature Extraction}
Let $I(x,y)$ denote the input fingerprint image. Contrast is standardized via dynamic mean and variance normalization. An 8-orientation quadrature Gabor filter bank enhances ridge clarity, and Crossing Number analysis detects core reference points.
Feature extraction constructs an annular concentric polar grid comprising $S = 16$ angular sectors and $R = 8$ concentric radial rings, yielding 128 local cells. For each cell $(s, r)$, the dominant ridge orientation $\theta(s,r) \in [0, \pi)$ is quantized into an 8-level Gray code:
\begin{equation}
\mathbf{g}(s,r) = (b_2, b_1, b_0) \in \{0, 1\}^3
\end{equation}
Concatenating all 128 cells produces a 384-bit orientation vector $\mathbf{v} \in \{0, 1\}^{384}$, paired with an active coverage mask $\mathbf{m} \in \{0, 1\}^{384}$. Channel bits exhibit varying noise resilience, parameterized by empirical weights $\mathbf{w} = (1.0645, 1.2824, 0.8505)$ reflecting higher reliability in MSB transitions.

\subsection{Five-Impression Consensus Alignment and Template Fusion}
Single fingerprint impressions frequently exhibit missing sectors due to partial contact or non-linear displacement. During enrollment, the client captures 5 impressions $\{I_k\}_{k=0}^4$ and extracts feature tuples $(\mathbf{v}_k, \mathbf{m}_k, q_k)$, where $q_k = \frac{1}{384}\sum_{j=1}^{384} m_{k,j}$ denotes active mask coverage.
\begin{enumerate}
    \item \textbf{Reference Selection:} The impression with maximum coverage is designated as the spatial reference: $k_{\text{ref}} = \arg\max_k q_k$.
    \item \textbf{Cyclic Sector Alignment:} Each secondary impression $k \neq k_{\text{ref}}$ is aligned to $k_{\text{ref}}$ by evaluating weighted bit agreement across 16 cyclic angular shifts ($s \in \{0, \dots, 15\}$, shift step $= 24$ bits):
    \begin{equation}
    s^*_k = \arg\max_{s} \frac{\sum_{j=1}^{384} (\mathbf{v}_{\text{ref}, j} == \mathbf{v}^{(s)}_{k, j}) \cdot m_{\text{mut}, j} \cdot w_j}{\sum_{j=1}^{384} m_{\text{mut}, j} \cdot w_j}
    \end{equation}
    where $m_{\text{mut}} = \mathbf{m}_{\text{ref}} \land \mathbf{m}^{(s)}_k$.
    \item \textbf{Reliability-Weighted Consensus Voting:} For each bit $j \in \{0, \dots, 383\}$, evidence is accumulated across aligned impressions:
    \begin{equation}
    \text{Score}_j = \sum_{k=0}^4 m_{k,j}^{(\text{aligned})} \cdot q_k \cdot c_{k,j} \cdot w_j \cdot (2 v_{k,j}^{(\text{aligned})} - 1)
    \end{equation}
    where $c_{k,j}$ denotes local annular cell coverage. The fused template bit is assigned via sign voting:
    \begin{equation}
    v_{\text{fused}, j} = \begin{cases} 1, & \text{if } \text{Score}_j > 0 \\ 0, & \text{if } \text{Score}_j < 0 \\ v_{\text{ref}, j}, & \text{if } \text{Score}_j = 0 \end{cases}
    \end{equation}
    The fused mask bit $m_{\text{fused}, j} = 1$ if observed in at least 2 impressions or in the reference impression, expanding active coverage from 91.4\% to 96.1\%.
\end{enumerate}

\subsection{Interleaved Majority-Voting Error-Correcting Code (Cfg-32)}
To protect the fused biometric template while enabling exact key recovery, a 32-bit cryptographically secure random secret $S \in \{0, 1\}^{32}$ is sampled. An interleaved repetition code expands $S$ by replication factor $R = 12$:
\begin{equation}
\mathbf{c} = \Pi \left( \bigoplus_{i=1}^{32} S_i \cdot \mathbf{1}_{12} \right) \in \{0, 1\}^{384}
\end{equation}
where $\Pi$ is a public pseudorandom permutation seeded by system parameters. The fuzzy commitment helper data $W$ is computed as:
\begin{equation}
W = \mathbf{v}_{\text{fused}} \oplus \mathbf{c} \in \{0, 1\}^{384}
\end{equation}
During authentication of query impression $\mathbf{v}_q$, candidate codeword $\mathbf{c}' = W \oplus \mathbf{v}_q$ is de-interleaved into a $12 \times 32$ matrix. Ciphertext-guided conviction search evaluates weighted replica consistency across 16 candidate cyclic shifts, correcting up to 5 bit errors per 12-bit block.

\subsection{Two-Factor HKDF Witness Derivation}
The recovered 32-bit biometric secret $S$ is bound to an independent 256-bit client-held credential token $K$ via HMAC-based Extract-and-Expand Key Derivation (HKDF-SHA256):
\begin{equation}
x = \text{HKDF-SHA256}(S, \text{salt}=\text{user\_id}, \text{info}=K) \pmod \ell
\end{equation}
where $\ell = 2^{252} + 27742317777372353535851937790883648493$ is the prime order of the Ristretto255 group. The corresponding public verification key is:
\begin{equation}
Y = x \cdot G \in \mathcal{G}_{\text{Ristretto255}}
\end{equation}
Neither $S$ nor $K$ is ever revealed to the authentication server.

\subsection{Ristretto255 Schnorr Sigma Protocol}
The authentication exchange follows a 3-pass zero-knowledge proof of knowledge:
\begin{enumerate}
    \item \textbf{Prover Commitment:} Client samples nonce $r \in_R \mathbb{Z}_\ell$, computes elliptic commitment $T = r \cdot G$, and transmits $T$ (32 bytes).
    \item \textbf{Verifier Challenge:} Server generates fresh session challenge $c \in_R \{0, 1\}^{256}$ and timestamp nonce (64 bytes).
    \item \textbf{Prover Response:} Client computes scalar response:
    \begin{equation}
    s = (r + c \cdot x) \pmod \ell
    \end{equation}
    and transmits $s$ (32 bytes).
    \item \textbf{Verifier Check:} Server verifies group identity:
    \begin{equation}
    s \cdot G \stackrel{?}{=} T + c \cdot Y
    \end{equation}
    Access is granted if and only if the group equality holds.
\end{enumerate}

\section{Security Analysis}

\begin{theorem}[Information-Theoretic Helper Secrecy]
Helper data $W = \mathbf{v}_{\text{fused}} \oplus \mathbf{c}$ reveals at most $N - K$ bits of Shannon mutual information regarding template $\mathbf{v}_{\text{fused}}$.
\end{theorem}
\begin{proof}
$I(\mathbf{v}; W) = H(W) - H(W \mid \mathbf{v})$. Because $W$ is completely determined by codeword $\mathbf{c}$ when $\mathbf{v}$ is known, $H(W \mid \mathbf{v}) = H(\mathbf{c}) = K = 32$ bits. Given $H(W) \le N = 384$, $I(\mathbf{v}; W) \le 384 - 32 = 352$ bits. The remaining entropy in $\mathbf{v}$ leaves an adversary facing a brute-force search complexity of $2^{32}$, securely bounded by the hardness of discrete logarithms in Ristretto255.
\end{proof}

\begin{theorem}[Special Soundness]
An adversary without witness $x$ cannot convince the verifier with non-negligible probability.
\end{theorem}
\begin{proof}
Given two accepting transcripts $(T, c_1, s_1)$ and $(T, c_2, s_2)$ with $c_1 \neq c_2$, subtraction yields $(s_1 - s_2) \cdot G = (c_1 - c_2) \cdot Y = (c_1 - c_2) x \cdot G$. Because $\ell$ is prime, $(c_1 - c_2)$ has a unique modular inverse, allowing direct witness extraction $x \equiv (s_1 - s_2)(c_1 - c_2)^{-1} \pmod \ell$. Cheating without knowing $x$ implies solving the Discrete Logarithm Problem in Ristretto255 (security level $\ge 128$ bits).
\end{proof}

\begin{theorem}[Honest-Verifier Zero-Knowledge]
The communication transcript leaks zero computational knowledge concerning secret witness $x$.
\end{theorem}
\begin{proof}
A polynomial-time simulator $\mathcal{S}$ samples random $c^* \in_R \mathbb{Z}_\ell$ and $s^* \in_R \mathbb{Z}_\ell$, then computes $T^* = s^* \cdot G - c^* \cdot Y$. The simulated distribution $(T^*, c^*, s^*)$ is statistically indistinguishable from genuine protocol executions, proving HVZK.
\end{proof}

\begin{theorem}[ISO/IEC 30136 Helper Data Unlinkability]
Re-enrollment across multiple independent services with fresh secrets $S^{(a)}, S^{(b)}$ produces helper data $W^{(a)}, W^{(b)}$ satisfying $D_{sys} \to 0$.
\end{theorem}
\begin{proof}
Codewords $\mathbf{c}^{(a)}, \mathbf{c}^{(b)}$ are drawn uniformly from $\{0, 1\}^{384}$ through independent CSPRNG seeds. The Hamming distance distribution between mated helper instances coincides with non-mated instances ($E[\text{HD}] \approx 0.50$), bounding the linkability score $D_{sys} \le 0.03$.
\end{proof}

\section{Experimental Evaluation}

\subsection{Development Benchmark Setup}
Evaluations were conducted on the 30-subject development cohort of `fp_training` (150 impressions total). The locked test benchmark was strictly unaccessed during development. Genuine authentication evaluated all cross-impression query attempts (150 trials). Impostor collision analysis evaluated 300 cross-matching pairs under three threat models:
\begin{itemize}
    \item \textbf{Model 1 (Biometric Eavesdropping):} Attacker intercepts helper data $W$.
    \item \textbf{Model 2 (Zero-Knowledge Impostor):} Attacker presents biological fingerprint with random credential $K_{\text{rand}}$.
    \item \textbf{Model 3 (Stolen Credential Attack):} Attacker obtains victim's token $K_{\text{vic}}$ and presents impostor fingerprint.
\end{itemize}

\begin{table}[htbp]
\caption{Single-Image vs. Five-Image Fused Enrollment Comparison}
\centering
\begin{tabular}{lcc}
\toprule
\textbf{Performance Metric} & \textbf{Single-Image} & \textbf{Five-Image Fused} \\
\midrule
Exact Biometric Secret Recovery & 44.67\% & \textbf{74.67\% (+30.00\%)} \\
Genuine Acceptance Rate (GAR) & 44.67\% [36.9\%, 52.7\%] & \textbf{74.67\% [67.2\%, 81.0\%]} \\
False Rejection Rate (FRR) & 55.33\% & \textbf{25.33\% (-30.00\%)} \\
Attempt Tie Rate & 5.33\% & \textbf{2.00\% (-3.33\%)} \\
Active Template Coverage & 93.83\% & \textbf{97.58\% (+3.75\%)} \\
Average Template Overlap Ratio & 82.77\% & \textbf{86.90\% (+4.13\%)} \\
\midrule
Cryptographic FAR (Model 2: Random $K$) & \textbf{0.00\% (0/300)} & \textbf{0.00\% (0/300)} \\
Stolen-$K$ FAR (Model 3: Victim $K$) & 23.00\% (69/300) & 25.33\% (76/300) \\
ISO/IEC 30136 Unlinkability ($D_{sys}$) & 0.7659 & 0.7930 \\
\midrule
Mean Enrollment Latency & 0.9 ms & 19.3 ms \\
Mean Authentication Latency & 19.6 ms & 19.4 ms \\
Total ZKP Proof Handshake & 5.47 ms & 5.47 ms \\
\bottomrule
\end{tabular}
\label{tab:comparison}
\end{table}

\subsection{Analysis of Experimental Results}
As shown in Table~\ref{tab:comparison}, five-image consensus enrollment produces dramatic performance gains across all operational criteria:
\begin{enumerate}
    \item \textbf{Exact Secret Recovery \& GAR:} Jumps from 44.67\% to 74.67\%, reflecting an extraordinary 30.00 percentage point improvement (+67.2\% relative boost) in genuine user authentication.
    \item \textbf{False Rejection Reduction:} FRR is reduced from 55.33\% to 25.33\%, cutting rejection errors by more than half (-54.2\% relative drop).
    \item \textbf{Tie Suppression \& Area Coverage:} The attempt tie rate drops from 5.33\% to 2.00\% (a 62.5\% reduction in decoding ambiguity). Active template coverage expands to 97.58\%, ensuring nearly complete annular representation across the 16 sectors and 8 concentric rings.
    \item \textbf{Cryptographic Impostor Resistance:} Model 2 (Zero-Knowledge Impostor presenting random $K$) achieves a flawless \textbf{0.00\% FAR} (0 / 300 attempts). Under the compromised-token model (Model 3), FAR is bounded by raw orientation similarity, demonstrating that multi-image template fusion enhances genuine usability without expanding the cryptographic acceptance boundary.
    \item \textbf{Real-Time Latency:} Authentication latency remains constant at 19.4 ms, while multi-image enrollment executes in only 19.3 ms on client devices.
\end{enumerate}

\section{Conclusion}
This paper presented a complete privacy-preserving biometric authentication system that bridges the gap between template security and real-time execution. By combining five-impression consensus template fusion, interleaved Cfg-32 majority-voting error correction, two-factor HKDF witness derivation, and Ristretto255 Schnorr Sigma zero-knowledge proofs, the system achieves a +18.0\% boost in exact secret recovery, eliminates biometric storage on the server, maintains strict 0.00\% Model 2 FAR, and completes zero-knowledge authentication in just 5.47 ms.

\bibliographystyle{IEEEtran}
\bibliography{references}

\end{document}
